\section{Topological rings and modules}
\label{section-topological-ring}

\noindent
Let's quickly discuss some properties of topological abelian groups.
An abelian group $M$ is a {\it topological abelian group} if $M$ is
endowed with a topology such that addition $M \times M \to M$,
$(x, y) \mapsto x + y$ and inverse $M \to M$, $x \mapsto -x$ are
continuous. A {\it homomorphism of topological abelian groups}
is just a homomorphism of abelian groups which is continuous.
The category of commutative topological groups is additive and
has kernels and cokernels, but is not abelian (as the axiom
$\Im = \Coim$ doesn't hold). If $N \subset M$ is a
subgroup, then we think of $N$ and $M/N$ as topological groups
also, namely using the induced topology on $N$ and the quotient
topology on $M/N$ (i.e., such that $M \to M/N$ is submersive).
Note that if $N \subset M$ is an open subgroup, then the topology
on $M/N$ is discrete.

\medskip\noindent
We say the topology on $M$ is {\it linear} if there exists a fundamental
system of neighbourhoods of $0$ consisting of subgroups. If so then these
subgroups are also open. An example is the following. Let $I$ be a directed set
and let $G_i$ be an inverse system of (discrete)
abelian groups over $I$. Then
$$
G = \lim_{i \in I} G_i
$$
with the inverse limit topology is linearly topologized with a fundamental
system of neighbourhoods of $0$ given by $\Ker(G \to G_i)$.
Conversely, let $M$ be a linearly topologized abelian group.
Choose any fundamental system of open subgroups $U_i \subset M$, $i \in I$
(i.e., the $U_i$ form a fundamental system of open neighbourhoods and
each $U_i$ is a subgroup of $M$). Setting
$i \geq i' \Leftrightarrow U_i \subset U_{i'}$ we see that $I$
is a directed set. We obtain a homomorphism of
linearly topologized abelian groups
$$
c : M \longrightarrow \lim_{i \in I} M/U_i.
$$
It is clear that $M$ is {\it separated} (as a topological space)
if and only if $c$ is injective. We say that $M$ is {\it complete}
if $c$ is an isomorphism\footnote{We include being separated as part
of being complete as we'd like to have a unique limits in complete
groups. There is a definition of completeness for any topological group,
agreeing, modulo the separation issue, with this one in our special case.}.
We leave it to the reader to check that this condition is independent of
the choice of fundamental system of open subgroups $\{U_i\}_{i \in I}$
chosen above. In fact the topological abelian group
$M^\wedge = \lim_{i \in I} M/U_i$ is independent of this
choice and is sometimes called the {\it completion} of $M$.
Any $G = \lim G_i$ as above is complete, in particular,
the completion $M^\wedge$ is always complete.

\begin{definition}[Topological rings]
\label{definition-topological-ring}
\begin{reference}
\cite[Chapter 0, Sections 7.1 and 7.2]{EGA1}
\end{reference}
Let $R$ be a ring and let $M$ be an $R$-module.
\begin{enumerate}
\item We say $R$ is a {\it topological ring} if $R$ is endowed with a topology
such that both addition and multiplication are continuous as maps
$R \times R \to R$ where $R \times R$ has the product topology.
In this case we say $M$ is a {\it topological module} if $M$ is endowed
with a topology such that addition $M \times M \to M$ and
scalar multiplication $R \times M \to M$ are continuous.
\item A {\it homomorphism of topological modules} is just a continuous
$R$-module map. A {\it homomorphism of topological rings} is a
ring homomorphism which is continuous for the given topologies.
\item We say $M$ is {\it linearly topologized} if $0$ has a fundamental
system of neighbourhoods consisting of submodules. We say $R$ is
{\it linearly topologized} if $0$ has a fundamental system of neighbourhoods
consisting of ideals.
\item If $R$ is linearly topologized, we say that $I \subset R$ is an
{\it ideal of definition} if $I$ is open and if every neighbourhood
of $0$ contains $I^n$ for some $n$.
\item If $R$ is linearly topologized, we say that $R$ is {\it pre-admissible}
if $R$ has an ideal of definition.
\item If $R$ is linearly topologized, we say that $R$ is {\it admissible} if
it is pre-admissible and
complete\footnote{By our conventions this includes separated.}.
\item If $R$ is linearly topologized, we say that $R$ is {\it pre-adic} if
there exists an ideal of definition $I$ such that $\{I^n\}_{n \geq 0}$
forms a fundamental system of neighbourhoods of $0$.
\item If $R$ is linearly topologized, we say that $R$ is {\it adic} if
$R$ is pre-adic and complete.
\end{enumerate}
Note that a (pre)adic topological ring is the same thing as a (pre)admissible
topological ring which has an ideal of definition $I$ such that $I^n$ is
open for all $n \geq 1$.
\end{definition}

\noindent
Let $R$ be a ring and let $M$ be an $R$-module. Let $I \subset R$ be
an ideal. Then we can consider the linear topology on $R$ which has
$\{I^n\}_{n \geq 0}$ as a fundamental system of neighbourhoods of $0$.
This topology is called the $I$-adic topology; $R$ is a pre-adic
topological ring in the $I$-adic topology\footnote{Thus the $I$-adic
topology is sometimes called the $I$-pre-adic topology.}. Moreover, the
linear topology
on $M$ which has $\{I^nM\}_{n \geq 0}$ as a fundamental system of open
neighbourhoods of $0$ turns $M$ into a topological $R$-module.
This is called the {\it $I$-adic topology} on $M$. We see that
$M$ is $I$-adically complete (as defined in
Algebra, Definition \ref{algebra-definition-complete})
if and only if $M$ is complete in the $I$-adic topology\footnote{
It may happen that the $I$-adic completion $M^\wedge$
is not $I$-adically complete, even though $M^\wedge$
is always complete with respect to the limit topology. If $I$ is finitely
generated then the $I$-adic topology and the limit topology on $M^\wedge$
agree, see Algebra, Lemma \ref{algebra-lemma-hathat-finitely-generated} and
its proof.}.
In particular, we see that $R$ is $I$-adically complete if and
only if $R$ is an adic topological ring in the $I$-adic topology.

\medskip\noindent
As a special case, note that the discrete topology is the $0$-adic
topology and that any ring in the discrete topology is adic.

\begin{lemma}
\label{lemma-continuous}
Let $\varphi : R \to S$ be a ring map.
Let $I \subset R$ and $J \subset S$ be ideals
and endow $R$ with the $I$-adic topology and $S$ with the $J$-adic
topology. Then $\varphi$ is a homomorphism of topological rings
if and only if $\varphi(I^n) \subset J$ for some $n \geq 1$.
\end{lemma}

\begin{proof}
Omitted.
\end{proof}

\begin{lemma}[Baire category theorem]
\label{lemma-baire-category-complete-module}
Let $M$ be a topological abelian group. Assume $M$ is linearly
topologized, complete, and has a countable fundamental system of
neighbourhoods of $0$. If $U_n \subset M$, $n \geq 1$
are open dense subsets, then $\bigcap_{n \geq 1} U_n$ is dense.
\end{lemma}

\begin{proof}
Let $U_n$ be as in the statement of the lemma. After replacing $U_n$ by
$U_1 \cap \ldots \cap U_n$, we may assume that
$U_1 \supset U_2 \supset \ldots$.
Let $M_n$, $n \in \mathbf{N}$ be a fundamental
system of neighbourhoods of $0$. We may assume that
$M_{n + 1} \subset M_n$. Pick $x \in M$. We will show that
for every $k \geq 1$ there exists a $y \in \bigcap_{n \geq 1} U_n$
with $x - y \in M_k$.

\medskip\noindent
To construct $y$ we argue as follows.
First, we pick a $y_1 \in U_1$ with $y_1 \in x + M_k$.
This is possible because $U_1$ is dense and $x + M_k$ is open.
Then we pick a $k_1 > k$ such that $y_1 + M_{k_1} \subset U_1$.
This is possible because $U_1$ is open.
Next, we pick a $y_2 \in U_2$ with $y_2 \in y_1 + M_{k_1}$.
This is possible because $U_2$ is dense and $y_2 + M_{k_1}$ is open.
Then we pick a $k_2 > k_1$ such that $y_2 + M_{k_2} \subset U_2$.
This is possible because $U_2$ is open.

\medskip\noindent
Continuing in this fashion we get a converging sequence
$y_i$ of elements of $M$ with limit $y$. By construction
$x - y \in M_k$. Since
$$
y - y_i = (y_{i + 1} - y_i) + (y_{i + 2} - y_{i + 1}) + \ldots
$$
is in $M_{k_i}$ we see that $y \in y_i + M_{k_i} \subset U_i$
for all $i$ as desired.
\end{proof}

\begin{lemma}
\label{lemma-consequence-baire-complete-module}
With same assumptions as Lemma \ref{lemma-baire-category-complete-module}
if $M = \bigcup_{n \geq 1} N_n$ for some closed subgroups $N_n$,
then $N_n$ is open for some $n$.
\end{lemma}

\begin{proof}
If not, then $U_n = M \setminus N_n$ is dense for all $n$ and
we get a contradiction with Lemma \ref{lemma-baire-category-complete-module}.
\end{proof}

\begin{lemma}[Open mapping lemma]
\label{lemma-open-mapping}
Let $u : N \to M$ be a continuous map of topologized
abelian groups. Assume that $M$ separated and that $N$ is complete, linearly
topologized, and has a countable fundamental system of neighbourhoods of $0$.
Then exactly one of the following holds
\begin{enumerate}
\item $u$ is open, or
\item for some open subgroup $N' \subset N$ the image
$u(N')$ is nowhere dense in $M$.
\end{enumerate}
\end{lemma}

\begin{proof}
Let $N_n$, $n \in \mathbf{N}$ be a fundamental system of neighbourhoods
of $0$. We may assume that $N_n$ is a subgroup and that
$N_{n + 1} \subset N_n$. If (2) does not hold, then
the closure $M_n$ of $u(N_n)$ is an open subgroup for $n = 1, 2, 3, \ldots$.
Since $u$ is continuous, we see that $M_n$, $n \in \mathbf{N}$
must be a fundamental system of open neighbourhoods of $0$ in $M$.
Also, since $M_n$ is the closure of $u(N_n)$ we see that
$$
u(N_n) + M_{n + 1} = M_n
$$
for all $n \geq 1$. Pick $x_1 \in M_1$. Then we can inductively choose
$y_i \in N_i$ and $x_{i + 1} \in M_{i + 1}$ such that
$$
u(y_i) + x_{i + 1} = x_i
$$
The element $y = y_1 + y_2 + y_3 + \ldots$ of $N$ exists because
$N$ is complete. Whereupon we see that $x_1 = u(y)$ because $M$ is separated.
Thus $M_1 = u(N_1)$. In exactly the same way the reader shows that
$M_i = u(N_i)$ for all $i \geq 2$ and we see that $u$ is open.
\end{proof}









